\section{中美格局：如何跨越地缘封锁}

上一章讨论公司图谱，本章收束到中美格局。访谈最后强调，科技投资和科技创业不能靠混圈子，还是要靠创造。地缘封锁会限制算力、资本、市场和合作，但不会自动决定胜负；真正穿越限制的，是技术创造力、工程执行、产品洞察、开源传播和产业链组织能力。

\begin{figure}[H]
\centering
\includegraphics[width=0.92\textwidth]{china-us-geofence.png}
\caption{中美 AI 格局：地缘封锁下，技术创造力比混圈更重要。自制概念图，依据 01:54:32--02:01:11 对谈内容整理。}
\end{figure}

\begin{knowledgebox}{读图：优势不同，不代表只能防守}
美国侧有算力、frontier lab、资本和全球产品入口，中国侧有工程速度、应用场景、开源效率和产业链执行力。地缘封锁会制造摩擦，但也会迫使团队在效率、工程和本地生态上找新路径。
\end{knowledgebox}

\subsection{中国团队的机会在哪里}

本节把机会拆成四类。第一是效率机会：用更少算力做出可竞争模型和产品。第二是应用机会：中国有大量数字化业务、内容、电商、工业和硬件场景。第三是开源机会：通过开放模型和工具获得全球开发者分发。第四是产业链机会：AI 与硬件、机器人、汽车、消费电子结合时，中国供应链和工程速度会变得重要。

\begin{importantbox}{实践经验：封锁时代更需要可验证创造}
越是外部约束强，越不能只讲叙事。团队需要拿出可验证的模型、代码、产品、收入、效率曲线和开发者采用。真正的全球影响力来自创造物，而不是圈层背书。
\end{importantbox}

\subsection{地缘约束下的风险}

本节同样要保留风险意识。中国团队面临算力供应、海外合规、品牌信任、企业销售、支付和生态入口等限制；美国团队也面临成本、监管、产品泡沫、资本预期和组织复杂化。中美格局不是单方面压制，也不是单方面反超，而是两种创新系统在不同约束下寻找突破。

\noindent\begin{tabularx}{\textwidth}{p{0.18\textwidth}p{0.34\textwidth}X}
\toprule
维度 & 美国优势/风险 & 中国优势/风险 \\
\midrule
算力 & 供应和云生态强，但成本和监管压力高 & 受出口管制约束，需要效率和替代方案。 \\
模型 & frontier lab 密集 & 开源效率和工程 recipe 可能成为突破口。 \\
产品 & 全球入口和企业客户强 & 本地场景丰富，但全球化品牌更难。 \\
产业链 & 软件生态强 & 硬件、汽车、机器人和制造链条强。 \\
\bottomrule
\end{tabularx}

\subsection{本章小结}

中美 AI 格局的核心，不是简单乐观或悲观，而是看谁能在约束下创造更好的模型、产品、环境和生态。对 AI/互联网方向的长期追踪，应关注可验证产物，而不是只关注叙事热度。

